%% ----------------------------------------------------------------------------
%%  §7.1  Funzioni di due variabili: dominio, grafico, curve di livello
%% ----------------------------------------------------------------------------
\section{Funzioni di due variabili: dominio, grafico, curve di livello}
\label{sec:due-variabili}

Indichiamo con $\R^{2}$ l'insieme delle coppie ordinate di numeri reali:
\begin{equation}\label{eq:R2}
  \R^{2}=\set{(x,y) : x\in\R,\ y\in\R}.
\end{equation}
Fissato nel piano un riferimento cartesiano ortogonale, ogni coppia $(x,y)$ individua un
punto $P$ del piano, e viceversa. La \emph{distanza} tra due punti $P_{1}=(x_{1},y_{1})$ e
$P_{2}=(x_{2},y_{2})$ è, per il teorema di Pitagora,
\begin{equation}\label{eq:distanza}
  d(P_{1},P_{2})=\sqrt{(x_{1}-x_{2})^{2}+(y_{1}-y_{2})^{2}} .
\end{equation}
Ci servirà spesso il confronto tra la distanza e le differenze delle coordinate: posto
$a=x_{1}-x_{2}$ e $b=y_{1}-y_{2}$, poiché $a^{2}\le a^{2}+b^{2}\le(\abs a+\abs b)^{2}$,
risulta
\begin{equation}\label{eq:distanza-componenti}
\begin{gathered}
  \abs{x_{1}-x_{2}}\le d(P_{1},P_{2}), \qquad \abs{y_{1}-y_{2}}\le d(P_{1},P_{2}),\\
  d(P_{1},P_{2})\le\abs{x_{1}-x_{2}}+\abs{y_{1}-y_{2}} .
\end{gathered}
\end{equation}
In parole: due punti sono vicini se e soltanto se sono vicine sia le ascisse sia le
ordinate.

\begin{definizione}[Funzione di due variabili]\label{def:funzione-due-variabili}
\index{funzione!di due variabili}\index{dominio}
Sia $D$ un sottoinsieme di $\R^{2}$. Una legge $f$ che a ogni elemento di $D$ fa
corrispondere uno e un solo numero reale si dice una \emph{funzione di due variabili}; si
scrive $f:D\to\R$, oppure, per mettere in evidenza che $D\subseteq\R^{2}$,
\begin{equation}\label{eq:z-f-xy}
  z=f(x,y), \qquad (x,y)\in D .
\end{equation}
L'insieme $D$ si dice il \emph{dominio} (o \emph{insieme di definizione}) di $f$.
L'insieme dei punti $(x,y,z)$ dello spazio con $(x,y)\in D$ e $z=f(x,y)$ si chiama il
\emph{grafico} di $f$: è, in generale, una superficie.
\end{definizione}

Per rappresentare graficamente una funzione di due variabili si usa un riferimento
cartesiano ortogonale di assi $x$, $y$, $z$: a ogni punto $(x,y)$ del dominio, nel piano
di base, si associa il punto $(x,y,z)$ a quota $z=f(x,y)$. Un ingegnere aeronautico
incontra funzioni di due variabili a ogni passo. La portanza di un velivolo al livello
del mare, per una data superficie alare $S$, è
\begin{equation}\label{eq:portanza-due-variabili}
  L(V,C_{L})=\frac12\rho S V^{2}C_{L},
\end{equation}
una funzione della velocità $V$ e del coefficiente di portanza $C_{L}$. Prima ancora dei
grafici, il modo tradizionale di rappresentare una funzione di due variabili è la
\emph{tabella a doppia entrata}, come quelle dell'atmosfera standard o dei prontuari.

\begin{python}[Una tabella a doppia entrata: la portanza $L(V,C_{L})$]
rho, S, W = 1.225, 16.0, 10000.0         # kg/m^3, m^2, N
L = lambda V, CL: 0.5 * rho * S * V**2 * CL

CLs = (0.2, 0.4, 0.6, 0.8, 1.0, 1.2)
print(" V (m/s) |" + "".join(f"   CL={c:.1f}" for c in CLs))
for V in (30, 40, 50, 60, 70):
    riga = "".join(f"{L(V, c):9.0f}" for c in CLs)
    print(f"  {V:5d}  |{riga}")
# le coppie (V, C_L) per cui L = W: C_L = 2W / (rho S V^2)
for V in (30, 40, 50, 60, 70):
    print(f"V = {V} m/s:  L = W per C_L = "
          f"{2 * W / (rho * S * V**2):.3f}")
\end{python}
\begin{risultato}
 V (m/s) |   CL=0.2   CL=0.4   CL=0.6   CL=0.8   CL=1.0   CL=1.2
     30  |     1764     3528     5292     7056     8820    10584
     40  |     3136     6272     9408    12544    15680    18816
     50  |     4900     9800    14700    19600    24500    29400
     60  |     7056    14112    21168    28224    35280    42336
     70  |     9604    19208    28812    38416    48020    57624
V = 30 m/s:  L = W per C_L = 1.134
V = 40 m/s:  L = W per C_L = 0.638
V = 50 m/s:  L = W per C_L = 0.408
V = 60 m/s:  L = W per C_L = 0.283
V = 70 m/s:  L = W per C_L = 0.208
\end{risultato}

Le ultime righe rispondono alla domanda che interessa al pilota: per sostenere il peso
$W=10$~kN, quale coefficiente di portanza serve a ciascuna velocità? Le coppie $(V,C_{L})$
per cui $L(V,C_{L})=W$ formano una curva del piano $(V,C_{L})$, che è un esempio della
nozione seguente.

\begin{definizione}[Curve di livello]\label{def:curve-livello}\index{curva di livello}
Sia $f:D\to\R$ e sia $c\in\R$. L'insieme
\begin{equation}\label{eq:curva-livello}
  \set{(x,y)\in D : f(x,y)=c}
\end{equation}
si chiama \emph{insieme di livello} (o, quando è una curva, \emph{curva di livello}) di
$f$ relativo al valore $c$.
\end{definizione}

La curva di livello $c$ è la proiezione sul piano di base della sezione del grafico con
il piano orizzontale $z=c$. Le carte topografiche rappresentano il terreno con le sue
curve di livello (le isoipse), le carte del tempo rappresentano la pressione con le
isobare: in entrambi i casi una funzione di due variabili è descritta, su un foglio
piano, dalla famiglia delle sue curve di livello (\cref{fig:livello}). Dove le curve di
livello sono fitte, la funzione varia rapidamente; dove sono rade, varia lentamente.

\begin{figure}[htb]
\centering
\includegraphics[width=\textwidth]{livello}
\caption{Due funzioni di due variabili descritte con le loro curve di livello. A
sinistra, le isobare (in hPa) di un campo di pressione con un'alta $A$ e una bassa $B$; le
frecce sono i vettori $-\nabla p$, di cui parleremo nel \cref{sec:derivate-parziali}. A
destra, le curve di livello della portanza $L(V,C_{L})$ di un velivolo con $S=16$~m$^{2}$
al livello del mare; la curva in rame è quella del volo livellato, $L=W$
(\texttt{python/fig\_livello.py}).}
\label{fig:livello}
\end{figure}

\subsection{Esempi di grafici}

Il comportamento di una funzione di due variabili si comprende spesso \emph{fissando una
delle due variabili} e studiando la funzione di una variabile che si ottiene.
Consideriamo
\begin{equation}\label{eq:sella}
  f(x,y)=x^{2}-y^{2},
\end{equation}
definita per ogni $(x,y)\in\R^{2}$. Per $y$ fissato si ottengono le parabole
\emph{convesse} $z=x^{2}-\text{costante}$, per $x$ fissato le parabole \emph{concave}
$z=\text{costante}-y^{2}$. Il grafico, rappresentato nella \cref{fig:superfici} sul
quadrato $[-1,1]\times[-1,1]$, ha la forma di una sella e si chiama \emph{paraboloide
iperbolico}: le sue curve di livello $x^{2}-y^{2}=c$ sono iperboli. Scambiando i ruoli
delle variabili si ottiene il grafico di $f(x,y)=y^{2}-x^{2}$, ruotato di un angolo retto.

La funzione
\begin{equation}\label{eq:rigata}
  f(x,y)=y\,(x^{2}+x)
\end{equation}
ha, per $y>0$ fissato, l'andamento di una parabola convessa del tipo $z=x^{2}+x$, mentre
per $y<0$ fissato l'andamento è quello di una parabola concava. Per $x$ fissato, invece,
la funzione è lineare: si tratta della retta $z=\text{costante}\cdot y$. Il grafico è
quindi unione di una famiglia di rette, e per questo si dice una \emph{superficie
rigata}.

\begin{figure}[htb]
\centering
\includegraphics[width=0.95\textwidth]{superfici}
\caption{I grafici di quattro funzioni di due variabili: il paraboloide iperbolico, la
superficie rigata $y(x^{2}+x)$, la funzione $\cos(x^{2}+y^{2})$, invariante per
rotazioni, e la funzione $\sen(xy)$, costante sulle iperboli $xy=t$
(\texttt{python/fig\_superfici.py}).}
\label{fig:superfici}
\end{figure}

La funzione
\begin{equation}\label{eq:cos-radiale}
  f(x,y)=\cos(x^{2}+y^{2})
\end{equation}
è definita per ogni $(x,y)\in\R^{2}$. Per comprenderne il grafico conviene pensarla come
funzione composta: posto $t=\sqrt{x^{2}+y^{2}}\ge0$, distanza del punto $(x,y)$
dall'origine, è $f(x,y)=\cos t^{2}$. La funzione è costante su ogni circonferenza
$x^{2}+y^{2}=t^{2}$ di centro l'origine, e il suo valore è $\cos t^{2}$. Il grafico si
ottiene quindi facendo ruotare intorno all'asse $z$ il profilo $z=\cos t^{2}$, $t\ge0$:
si dice che $f$ è \emph{invariante per rotazioni}. Più in generale, ogni funzione della
forma $f(x,y)=g\bigl(\sqrt{x^{2}+y^{2}}\bigr)$ è invariante per rotazioni, e il suo
grafico è una superficie di rotazione. La funzione
\begin{equation}\label{eq:sen-xy}
  f(x,y)=\sen(xy)
\end{equation}
è invece costante sulle iperboli $xy=t$, con $t$ fissato in $\R$.

Sono invarianti per rotazioni anche le funzioni
\begin{equation}\label{eq:radiali}
\begin{gathered}
  \log(x^{2}+y^{2}), \qquad f_{0}(x,y)=1-\sqrt{x^{2}+y^{2}},\\
  f_{1}(x,y)=1-\sqrt{x^{2}+y^{2}-1}, \qquad f_{2}(x,y)=1-\sqrt{x^{2}+y^{2}-2},
\end{gathered}
\end{equation}
i cui grafici sono nella \cref{fig:radiali}. La prima non è definita nell'origine, dove
$x^{2}+y^{2}=0$; il suo profilo è $z=\log t^{2}=2\log t$, $t>0$. La funzione $f_{0}$,
definita su tutto $\R^{2}$, ha per $y=0$ il profilo $z=1-\sqrt{x^{2}}=1-\abs x$: il grafico è
un cono, con il vertice appuntito nel punto $(0,0,1)$, dovuto al valore assoluto. La funzione
$f_{1}$ è definita se $x^{2}+y^{2}\ge1$, cioè all'esterno del cerchio di centro
l'origine e raggio $1$ (circonferenza compresa); analogamente $f_{2}$ è definita
all'esterno del cerchio di raggio $\sqrt2$.

\begin{figure}[htb]
\centering
\includegraphics[width=0.95\textwidth]{radiali}
\caption{Quattro funzioni invarianti per rotazioni: i grafici sono superfici di
rotazione intorno all'asse $z$. Le ultime due sono definite soltanto all'esterno di un
cerchio, sulla cui circonferenza il grafico ha tangente verticale
(\texttt{python/fig\_radiali.py}).}
\label{fig:radiali}
\end{figure}

\subsection{La determinazione del dominio}

Quando una funzione è assegnata mediante un'espressione analitica, senza specificare il
dominio, si sottintende che il dominio sia l'insieme di tutti i punti in cui
l'espressione ha significato. Determinarlo richiede di risolvere disequazioni in due
variabili, la cui soluzione è una regione del piano. Ad esempio, le funzioni
\begin{equation}\label{eq:domini-esempi}
\begin{gathered}
  z=1-\sqrt{(1-x^{2})(1-y^{2})},\\
  z=\sqrt{(x^{2}+y^{2}-x)(2x-x^{2}-y^{2})}
\end{gathered}
\end{equation}
non sono definite su tutto $\R^{2}$. La prima è definita quando
$(1-x^{2})(1-y^{2})\ge0$, cioè quando i fattori $1-x^{2}$ e $1-y^{2}$ hanno lo stesso
segno (o uno dei due è nullo): o entrambi $\ge0$, cioè $\abs x\le1$ e $\abs y\le1$ (il
quadrato centrale), oppure entrambi $\le0$, cioè $\abs x\ge1$ e $\abs y\ge1$ (le quattro
regioni d'angolo); si veda la \cref{fig:domini}, a sinistra.

La seconda è definita quando l'argomento della radice è $\ge0$. Completando i quadrati,
\begin{equation}\label{eq:due-cerchi}
\begin{gathered}
  x^{2}+y^{2}-x=\Bigl(x-\frac12\Bigr)^{2}+y^{2}-\frac14,\\
  2x-x^{2}-y^{2}=1-\bigl[(x-1)^{2}+y^{2}\bigr];
\end{gathered}
\end{equation}
il primo fattore è $\ge0$ all'esterno della circonferenza $C_{1}$ di centro
$(\frac12,0)$ e raggio $\frac12$, il secondo è $\ge0$ all'interno della circonferenza
$C_{2}$ di centro $(1,0)$ e raggio $1$ (circonferenze comprese). Il prodotto è $\ge0$
quando i fattori hanno lo stesso segno. Il caso «entrambi negativi» è impossibile: un
punto interno a $C_{1}$ è anche interno a $C_{2}$, perché $C_{1}$ è contenuta nel cerchio
delimitato da $C_{2}$ (le due circonferenze sono tangenti internamente nell'origine).
Resta il caso «entrambi $\ge0$»: il dominio è la lunula compresa tra le due
circonferenze (\cref{fig:domini}, a destra).

\begin{figure}[htb]
\centering
\begin{tikzpicture}[scale=0.85]
  \begin{scope}
    \clip (-2.3,-2.3) rectangle (2.3,2.3);
    \fill[pattern=north east lines, pattern color=blu] (-1,-1) rectangle (1,1);
    \fill[pattern=north east lines, pattern color=blu] (1,1) rectangle (2.3,2.3);
    \fill[pattern=north east lines, pattern color=blu] (-1,1) rectangle (-2.3,2.3);
    \fill[pattern=north east lines, pattern color=blu] (1,-1) rectangle (2.3,-2.3);
    \fill[pattern=north east lines, pattern color=blu] (-1,-1) rectangle (-2.3,-2.3);
    \draw[blu, thick] (-1,-2.3) -- (-1,2.3) (1,-2.3) -- (1,2.3)
      (-2.3,-1) -- (2.3,-1) (-2.3,1) -- (2.3,1);
  \end{scope}
  \draw[-Latex, grigio] (-2.5,0) -- (2.6,0) node[right, black] {$x$};
  \draw[-Latex, grigio] (0,-2.5) -- (0,2.6) node[above, black] {$y$};
  \node[font=\scriptsize, below left] at (-1,0) {$-1$};
  \node[font=\scriptsize, below right] at (1,0) {$1$};
  \node[font=\scriptsize, above left] at (0,1) {$1$};
  \node[font=\scriptsize, below left] at (0,-1) {$-1$};
\end{tikzpicture}\hspace{1cm}
\begin{tikzpicture}[scale=1.2]
  \fill[pattern=north east lines, pattern color=blu] (1,0) circle (1);
  \fill[white] (0.5,0) circle (0.5);
  \draw[blu, thick] (1,0) circle (1);
  \draw[blu, thick] (0.5,0) circle (0.5);
  \draw[-Latex, grigio] (-0.4,0) -- (2.4,0) node[right, black] {$x$};
  \draw[-Latex, grigio] (0,-1.3) -- (0,1.35) node[above, black] {$y$};
  \node[font=\scriptsize, below] at (0.5,0) {$\frac12$};
  \node[font=\scriptsize, below right] at (2,0) {$2$};
  \node[font=\scriptsize, blu] at (0.5,0.27) {$C_{1}$};
  \node[font=\scriptsize, blu] at (1.85,0.85) {$C_{2}$};
\end{tikzpicture}
\caption{I domini delle funzioni \eqref{eq:domini-esempi}. A sinistra,
$(1-x^{2})(1-y^{2})\ge0$: il quadrato centrale e le quattro regioni d'angolo. A destra,
la lunula compresa tra la circonferenza $C_{1}$ di centro $(\frac12,0)$ e raggio
$\frac12$ e la circonferenza $C_{2}$ di centro $(1,0)$ e raggio $1$.}
\label{fig:domini}
\end{figure}

\begin{notastorica}[Le curve di livello: da Halley alle carte del tempo]
La rappresentazione di una grandezza che varia sulla superficie terrestre mediante linee
che uniscono i punti in cui essa ha lo stesso valore risale all'astronomo Edmond Halley,
che nel 1701 pubblicò una carta dell'Atlantico con le linee di uguale declinazione
magnetica, preziose per la navigazione con la bussola. Nel 1752 il geografo francese
Philippe Buache pubblicò una carta del canale della Manica con le linee di uguale
profondità, e alla
fine del Settecento le curve di livello del terreno entrarono nella cartografia
topografica. Nell'Ottocento l'idea passò alla meteorologia: le carte sinottiche con le
isobare permisero di riconoscere i cicloni e gli anticicloni (questo termine fu coniato
da Francis Galton nel 1863). Oggi la stessa idea si ritrova nelle carte delle
prestazioni dei velivoli e dei motori, e nelle visualizzazioni dei risultati delle
simulazioni fluidodinamiche (CFD), dove le linee di uguale pressione o di uguale velocità
si chiamano \emph{isolinee}.
\end{notastorica}
